%
% Copyright (C) 1995, 2000, 2013 American Mathematical Society.
% Copyright (C) 2016-2025 LaTeX Project and American Mathematical Society.
%
% This work may be distributed and/or modified under the
% conditions of the LaTeX Project Public License, either version 1.3c
% of this license or (at your option) any later version.
% The latest version of this license is in
%   https://www.latex-project.org/lppl.txt
% and version 1.3c or later is part of all distributions of LaTeX
% version 2005/12/01 or later.
%
% This work has the LPPL maintenance status `maintained'.
%
% The Current Maintainer of this work is the LaTeX Project.
%
\documentclass[leqno,titlepage,openany]{amsldoc}[1999/12/13]

\def\MaintainedByLaTeXTeam#1{%
\begin{center}%
\fbox{\fbox{\begin{tabular}{@{}l@{}}%
This file is maintained by the \LaTeX{} Project team.\\%
Bug reports can be opened (category \texttt{#1}) at\\%
\url{https://latex-project.org/bugs/}.\end{tabular}}}\end{center}}

\ifx\UndEfiNed\url
  \ClassError{amsldoc}{%
    This version of amsldoc.tex must be processed\MessageBreak
    with a newer version of amsldoc.cls (2.02 or later).}{}
\fi

\renewcommand\bibname{References}

%    Definitions that should be moved to amsldoc.cls
%    when it is next updated
\makeatletter
\let\cleardouble@page\cleardoublepage
\AtBeginDocument{%
  \ifx\cleardouble@page\cleardoublepage
    \def\cleardoublepage{\clearpage{\pagestyle{empty}\cleardouble@page}}
  \fi
}

% #1 is empty or one of class, environment, option, package
% could use different font in each case
\makeatletter
\def\@category@index#1#2{%
    \autoindex{%
        #2\actualchar\string\index#1font{#2}%
        \ifx\@nil#1\@nil\else\space#1\fi
    }%
}

\let\indexpackagefont\textsf
\let\indexclassfont\textsf
\let\indexoptionfont\texttt
\let\indexenvironmentfont\texttt
\let\indexfont\texttt
\makeatother
%    For a package that won't be indexed, i.e. amsmath
\newcommand{\nipkg}{\textsf}
\def\tikz/{Ti\textit{k}Z}
%    End definitions for amsldoc.cls

\title{User's Guide for the \nipkg{amsmath} Package (Version~2.1)}
\author{American Mathematical Society, \LaTeX\ Project}
\date{1999-12-13\\(revised 2002-02-25, 2016-11-14, 2018-04-05, 2019-10-14, 2020-02-18)}
\makeatletter
\def\@thanks{\bigskip\MaintainedByLaTeXTeam{amslatex}}
\makeatother

% Very minor changes are not listed in the \date above:
%
% 2021-08-22: reported typo fixed mathools -> mathtools
% 2025-05-23: drop spurious option list


%    Use the amsmath package and amscd package in order to print
%    examples.
\usepackage{amsmath}
\usepackage{amscd}

\usepackage{imakeidx}
%    Need alternate coding for vertical bar to get it into index.
%    Based on \qcbang and \cnbang in amsldoc.cls, refined by David Carlisle.
%    Give it a sort field that will be at the end of the symbols;
%    must avoid " and | which have special meanings.
\DeclareRobustCommand{\qcvert}{\qc\|\index{{z1}@{\ntt\char'174}}}
\DeclareRobustCommand{\cnvert}{%
  \ncn{\|}\index{{z0}@{\ntt\bslchar\char'174}}}
%    Need alternate index sort field for bang to allow & to sort first
\DeclareRobustCommand{\altcnbang}{%
  \ncn{\!}\index{"!b@{\ntt\bslchar\qcbang}}}
\DeclareRobustCommand{\cnamp}{%
  \index{"!a@\texttt{\&}}}

\makeindex[intoc] % generate index data
\providecommand{\see}[2]{\textit{see} #1}

%    The amsldoc class includes a number of features useful for
%    documentation about TeX, including:
%
%    ---Commands \tex/, \amstex/, \latex/, ... for uniform treatment
%    of the various logos and easy handling of following spaces.
%
%    ---Commands for printing various common elements: \cn for command
%    names, \fn for file names (including font-file names), \env for
%    environments, \pkg and \cls for packages and classes, etc.

%    Many of the command names used here are rather long and will
%    contribute to poor linebreaking if we follow the \latex/ practice
%    of not hyphenating anything set in tt font; instead we selectively
%    allow some hyphenation.
\allowtthyphens % defined in amsldoc.cls

\hyphenation{ac-cent-ed-sym-bol add-to-counter add-to-length align-at
  aligned-at allow-dis-play-breaks ams-art ams-cd ams-la-tex amsl-doc
  ams-symb ams-tex ams-text ams-xtra bmatrix bold-sym-bol cen-ter-tags
  eqn-ar-ray idots-int int-lim-its latex med-space neg-med-space
  neg-thick-space neg-thin-space no-int-lim-its no-name-lim-its
  over-left-arrow over-left-right-arrow over-right-arrow pmatrix
  qed-sym-bol set-length side-set small-er tbinom the-equa-tion
  thick-space thin-space un-der-left-arrow un-der-left-right-arrow
  un-der-right-arrow use-pack-age var-inj-lim var-proj-lim vmatrix
  xalign-at xx-align-at}

%    Prepare for illustrating the \vec example
\newcommand{\vect}[1]{\mathbf{#1}}

\newcommand{\booktitle}[1]{\textit{#1}}
\newcommand{\journalname}[1]{\textit{#1}}
\newcommand{\seriesname}[1]{\textit{#1}}

%    Command to insert and index a particular phrase. Doesn't work for
%    certain kinds of special characters in the argument.
\newcommand{\ii}[1]{#1\index{#1}}

\newcommand{\vstrut}[1]{\vrule width0pt height#1\relax}

%    An environment for presenting comprehensive address information:
\newenvironment{infoaddress}{%
  \par\topsep\medskipamount
  \trivlist\centering
  \item[]%
  \begin{minipage}{.7\columnwidth}%
  \raggedright
}{%
  \end{minipage}%
  \endtrivlist
}

\newenvironment{eqxample}{%
  \par\addvspace\medskipamount
  \noindent\begin{minipage}{.5\columnwidth}%
  \def\producing{\end{minipage}\begin{minipage}{.5\columnwidth}%
    \hbox\bgroup\kern-.2pt\vrule width.2pt%
      \vbox\bgroup\parindent0pt\relax
%    The 3pt is to cancel the -\lineskip from \displ@y
    \abovedisplayskip3pt \abovedisplayshortskip\abovedisplayskip
    \belowdisplayskip0pt \belowdisplayshortskip\belowdisplayskip
    \noindent}
}{%
  \par
%    Ensure that a lonely \[\] structure doesn't take up width less than
%    \hsize.
  \hrule height0pt width\hsize
  \egroup\vrule width.2pt\kern-.2pt\egroup
  \end{minipage}%
  \par\addvspace\medskipamount
}

%    Errors in output shouldn't be \texttt, and they may also be longer
%    than one line, but they should be identified by the "errorbullet".
%    (They were formerly labeled as aubsections.)
\newenvironment{erroro}[1]{%
  \par\addvspace\medskipamount
  \noindent\hangindent\parindent
  \hbox to\parindent{\errorbullet\hfil}\ignorespaces
  #1\par\vspace{\smallskipamount}
}{%
  \par\addvspace\medskipamount
}

%    The chapters are so short, perhaps we shouldn't call them by the
%    name `Chapter'. We make \chaptername read an argument in order to
%    remove a following \space or "{} " (both possibilities are present
%    in book.cls).

\renewcommand{\chaptername}[1]{}
\newcommand{\chapnum}[1]{\mdash #1\mdash }
\makeatletter
\def\@makechapterhead#1{%
  \vspace{1.5\baselineskip}%
  {\parindent \z@ \raggedright \reset@font
    \ifnum \c@secnumdepth >\m@ne
      \large\bfseries \chapnum{\thechapter}%
      \par\nobreak
      \vskip.5\baselineskip\relax
    \fi
    #1\par\nobreak
    \vskip\baselineskip
  }}
\makeatother

%    A command for ragged-right parbox in a tabular.
\newcommand{\rp}{\let\PBS\\\raggedright\let\\\PBS}

%    Non-indexed file name
\newcommand{\nfn}[1]{\texttt{#1}}

%    For the examples in the math spacing table.
%%\newcommand{\lspx}{\mbox{\rule{5pt}{.6pt}\rule{.6pt}{6pt}}}
%%\newcommand{\rspx}{\mbox{\rule[-1pt]{.6pt}{7pt}%
%%  \rule[-1pt]{5pt}{.6pt}}}
\newcommand{\lspx}{\mathord{\Rightarrow\mkern-1mu}}
\newcommand{\rspx}{\mathord{\mkern-1mu\Leftarrow}}
\newcommand{\spx}[1]{$\lspx #1\rspx$}

%    For a list of characters representing document input.
\newcommand{\clist}[1]{%
  \mbox{\ntt\spaceskip.2em plus.1em \xspaceskip\spaceskip#1}}

%    Fix weird \latex/ definition of rightmark.
\makeatletter
\def\rightmark{\expandafter\@rightmark\botmark{}{}}
%    Also turn off section marks.
\let\sectionmark\@gobble
\renewcommand{\chaptermark}[1]{%
  \uppercase{\markboth{\rhcn#1}{\rhcn#1}}}
\newcommand{\rhcn}{\thechapter. }
\makeatother

%    Include down to \section but not \subsection, in toc:
\setcounter{tocdepth}{1}

%    Subheadings for the bibliography
\newif\iffirstbibsubhead \firstbibsubheadtrue
\newcommand{\bibsubhead}[1]{%
  \iffirstbibsubhead \firstbibsubheadfalse
  \else \addvspace{\medskipamount}
  \fi
  \item[]\hspace*{-\leftmargin}\textbf{#1}\par
  \vspace{\smallskipamount}
}

\DeclareMathOperator{\ix}{ix}
\DeclareMathOperator{\nul}{nul}
\DeclareMathOperator{\End}{End}
\DeclareMathOperator{\xxx}{xxx}

\usepackage{hyperref}
\hypersetup{pdfborder=0 0 0}
\makeatletter
\let\oldcs\cs
\def\cs#1{\texorpdfstring{\oldcs{#1}}{\@backslashchar\@backslashchar#1}}
\let\cn\cs
\makeatother
\begin{document}
\chapter{Displayed equations}

\section{Introduction}
The \nipkg{amsmath} package provides a number of additional displayed
equation structures\index{displayed equations}%
\index{equations|see{displayed equations}} beyond the ones
provided in basic \latex/. The augmented set includes:
\begin{verbatim}
  equation     equation*     align       align*
  gather       gather*       alignat     alignat*
  multline     multline*     flalign     flalign*
  split
\end{verbatim}
(Although the standard \env{eqnarray} environment remains available,
it is better to use \env{align} or \env{equation}+\env{split} instead.
Within \env{eqnarray}, spacing\index{horizontal spacing} around signs
of relation is not the preferred mathematical spacing, and is
inconsistent with that spacing as it appears in other environments.
Long lines in this environment may result in misplaced or overprinted
equation numbers. This environment also does not support the use of
\cn{qed} or \cn{qedhere} as provided by theorem packages.)

Except for \env{split}, each environment has both starred and unstarred
forms, where the unstarred forms have automatic numbering using
\latex/'s \env{equation} counter. You can suppress the number on any
particular line by putting \cn{notag} before the end of that line;
\cn{notag} should not be used outside a display environment as it will
mess up the numbering.  You can also
override\index{equationn@equation numbers!overriding} a number with a tag
of your own using \cn{tag}|{|\<label>|}|, where \<label> means arbitrary
text such as |$*$| or |ii| used to \qq{number} the equation.
A tag can reference a different tagged display by use of
|\tag{\ref{|\<label>|}|\<modifier>|}| where \<modifier> is optional.
If you are using \pkg{hyperref}, use \cn{ref*};
use of the starred form of \cn{ref} prevents a reference to a modified
tag containing a nested link from linking to the original display.

There is also a
\cn{tag*} command that causes the text you supply to be typeset
literally, without adding parentheses around it. \cn{tag} and \cn{tag*}
can also be used within the unnumbered versions of all the \nipkg{amsmath}
alignment structures. Some examples of the use of \cn{tag} may be found
in the sample files \fn{testmath.tex} and \fn{subeqn.tex} provided with
the \nipkg{amsmath} package.

The \env{split} environment is a special subordinate form that is used
only \emph{inside} one of the others. It cannot be used inside
\env{multline}, however. \env{split} supports only one alignment (|&|)
column; if more are needed, \env{aligned} or \env{alignedat} should be
used. The width of a \env{split} structure is the full line width.

In the structures that do alignment (\env{split}, \env{align} and
variants), relation symbols have an |&| before them but not
after\mdash unlike \env{eqnarray}. Putting the |&| after the
relation symbol will interfere with the normal spacing; it has to go
before.

In all multiline environments, lines are divided by \cn{\\}.  The \cn{\\}
should \emph{not} be used to end the last line.  Using it there will
result in unwanted extra vertical space following the display.

In \emph{all} math environments (inline or display), blank lines
(equivalent to \cn{par}) are not permitted, and will result in an error.

\begin{table}[p]
\caption[]{Comparison of displayed equation environments (vertical lines
indicating nominal margins)}\label{displays}
\renewcommand{\theequation}{\arabic{equation}}
\begin{eqxample}
\begin{verbatim}
\begin{equation*}
a=b
\end{equation*}
\end{verbatim}
\producing
\begin{equation*}
a=b
\end{equation*}
\end{eqxample}

\begin{eqxample}
\begin{verbatim}
\begin{equation}
a=b
\end{equation}
\end{verbatim}
\producing
\begin{equation}
a=b
\end{equation}
\end{eqxample}

\begin{eqxample}
\begin{verbatim}
\begin{equation}\label{xx}
\begin{split}
a& =b+c-d\\
 & \quad +e-f\\
 & =g+h\\
 & =i
\end{split}
\end{equation}
\end{verbatim}
\producing
\begin{equation}\label{xx}
\begin{split}
a& =b+c-d\\
 & \quad +e-f\\
 & =g+h\\
 & =i
\end{split}
\end{equation}
\end{eqxample}

\begin{eqxample}
\begin{verbatim}
\begin{multline}
a+b+c+d+e+f\\
+i+j+k+l+m+n
\end{multline}
\end{verbatim}
\producing
\begin{multline}
a+b+c+d+e+f\\
+i+j+k+l+m+n
\end{multline}
\end{eqxample}

\begin{eqxample}
\begin{verbatim}
\begin{gather}
a_1=b_1+c_1\\
a_2=b_2+c_2-d_2+e_2
\end{gather}
\end{verbatim}
\producing
\begin{gather}
a_1=b_1+c_1\\
a_2=b_2+c_2-d_2+e_2
\end{gather}
\end{eqxample}

\begin{eqxample}
\begin{verbatim}
\begin{align}
a_1& =b_1+c_1\\
a_2& =b_2+c_2-d_2+e_2
\end{align}
\end{verbatim}
\producing
\begin{align}
a_1& =b_1+c_1\\
a_2& =b_2+c_2-d_2+e_2
\end{align}
\end{eqxample}

\begin{eqxample}
\begin{verbatim}
\begin{align}
a_{11}& =b_{11}&
  a_{12}& =b_{12}\\
a_{21}& =b_{21}&
  a_{22}& =b_{22}+c_{22}
\end{align}
\end{verbatim}
\producing
\begin{align}
a_{11}& =b_{11}&
  a_{12}& =b_{12}\\
a_{21}& =b_{21}&
  a_{22}& =b_{22}+c_{22}
\end{align}
\end{eqxample}

\begin{eqxample}
\begin{verbatim}
\begin{flalign*}
a_{11} + b_{11}& = c_{11}&
  a_{12}& =b_{12}\\
b_{21}& = c_{21}&
  a_{22}& =b_{22}+c_{22}
\end{flalign*}
\end{verbatim}
\producing
\begin{flalign*}
a_{11} + b_{11}& = c_{11}&
  a_{12}& =b_{12}\\
b_{21}& = c_{21}&
  a_{22}& =b_{22}+c_{22}
\end{flalign*}
\end{eqxample}
\end{table}

%% ------------------------------------------------------------------ %%
\enlargethispage{1\baselineskip}

\section{Single equations}

The \env{equation} environment is for a single equation with an
automatically generated number. The \env{equation*} environment is the
same except for omitting the number.%
%%%%%%%%%%%%%%%%%%%%%%%%%%%%%%%%%%%%%%%%%%%%%%%%%%%%%%%%%%%%%%%%%%%%%%%%
\footnote{Basic \latex/ doesn't provide an \env{equation*} environment,
but rather a functionally equivalent environment named
\env{displaymath}.}
The wrapper \verb+\[ ... \]+ is equivalent to \env{equation*}.%
\index{[@\verb+\[ ... \]+}

%% ------------------------------------------------------------------ %%

\section{Split equations without alignment}

The \env{multline} environment is a variation of the \env{equation}
environment used for equations that don't fit on a single line. The
first line of a \env{multline} will be at the left margin and the last
line at the right margin, except for an indention on both sides in the
amount of \cn{multlinegap}. Any additional lines in between will be
centered independently within the display width (unless the \opt{fleqn}
option is in effect).%
\index{options!behavior of particular options}

Like \env{equation}, \env{multline} has only a single equation number
(thus, none of the individual lines should be marked with \cn{notag}).
The equation number is placed on the last line (\opt{reqno} option) or
first line (\opt{leqno} option); vertical centering as for \env{split}
is not supported by \env{multline}.%
\index{options!behavior of particular options}

It's possible to force one of the middle lines to the left or right with
commands \cn{shoveleft}, \cn{shoveright}. These commands take the entire
line as an argument, up to but not including the final \cn{\\}; for
example
\begin{multline}
\framebox[.65\columnwidth]{A}\\
\framebox[.5\columnwidth]{B}\\
\shoveright{\framebox[.55\columnwidth]{C}}\\
\framebox[.65\columnwidth]{D}
\end{multline}
\begin{verbatim}
\begin{multline}
\framebox[.65\columnwidth]{A}\\
\framebox[.5\columnwidth]{B}\\
\shoveright{\framebox[.55\columnwidth]{C}}\\
\framebox[.65\columnwidth]{D}
\end{multline}
\end{verbatim}

The value of \cn{multlinegap} can be changed with the usual \latex/
commands \cn{setlength} or \cn{addtolength}.

%% ------------------------------------------------------------------ %%

\section{Split equations with alignment}

Like \env{multline}, the \env{split} environment is for \emph{single}
equations that are too long to fit on one line and hence must be split
into multiple lines.  Unlike \env{multline}, however, the \env{split}
environment provides for alignment among the split lines, using |&| to
mark alignment points. Unlike the other \nipkg{amsmath} equation
structures, the \env{split} environment provides no numbering, because
it is intended to be used \emph{only inside some other displayed
  equation structure}, usually an \env{equation}, \env{align}, or
\env{gather} environment, which provides the numbering. For example:
\begin{equation}\label{e:barwq}\begin{split}
H_c&=\frac{1}{2n} \sum^n_{l=0}(-1)^{l}(n-{l})^{p-2}
  \sum_{l _1+\dots+ l _p=l}\prod^p_{i=1} \binom{n_i}{l _i}\\
&\quad\cdot[(n-l )-(n_i-l _i)]^{n_i-l _i}\cdot
  \Bigl[(n-l )^2-\sum^p_{j=1}(n_i-l _i)^2\Bigr].
  \kern-2em % adjust equation body to the right [mjd,13-Nov-1994]
\end{split}\end{equation}

\begin{verbatim}
\begin{equation}\label{e:barwq}\begin{split}
H_c&=\frac{1}{2n} \sum^n_{l=0}(-1)^{l}(n-{l})^{p-2}
  \sum_{l _1+\dots+ l _p=l}\prod^p_{i=1} \binom{n_i}{l _i}\\
&\quad\cdot[(n-l )-(n_i-l _i)]^{n_i-l _i}\cdot
  \Bigl[(n-l )^2-\sum^p_{j=1}(n_i-l _i)^2\Bigr].
\end{split}\end{equation}
\end{verbatim}

The \env{split} structure should constitute the entire body of the
enclosing structure, apart from commands like \cn{label} that produce no
visible material.

%% ------------------------------------------------------------------ %%

\section{Equation groups without alignment}

The \env{gather} environment is used for a group of consecutive
equations when there is no alignment desired among them; each one is
centered separately within the text width (see Table~\ref{displays}).
Equations inside \env{gather} are separated by a \cn{\\} command.
Any equation in a \env{gather} may consist of a \verb'\begin{split}'
  \dots\ \verb'\end{split}' structure\mdash for example:
\begin{verbatim}
\begin{gather}
  first equation\\
  \begin{split}
    second & equation\\
           & on two lines
  \end{split}
  \\
  third equation
\end{gather}
\end{verbatim}

%% ------------------------------------------------------------------ %%

\section{Equation groups with mutual alignment}

The \env{align} environment is used for two or more equations when
vertical alignment is desired; usually binary relations such as equal
signs are aligned (see Table~\ref{displays}).

To have several equation columns side-by-side, use extra ampersands
to separate the columns:
\begin{align}
x&=y       & X&=Y       & a&=b+c\\
x'&=y'     & X'&=Y'     & a'&=b\\
x+x'&=y+y' & X+X'&=Y+Y' & a'b&=c'b
\end{align}
%
\begin{verbatim}
\begin{align}
x&=y       & X&=Y       & a&=b+c\\
x'&=y'     & X'&=Y'     & a'&=b\\
x+x'&=y+y' & X+X'&=Y+Y' & a'b&=c'b
\end{align}
\end{verbatim}
Line-by-line annotations on an equation can be done by judicious
application of \cn{text} inside an \env{align} environment:
\begin{align}
x& = y_1-y_2+y_3-y_5+y_8-\dots
                    && \text{by \eqref{eq:C}}\\
 & = y'\circ y^*    && \text{by \eqref{eq:D}}\\
 & = y(0) y'        && \text {by Axiom 1.}
\end{align}
%
\begin{verbatim}
\begin{align}
x& = y_1-y_2+y_3-y_5+y_8-\dots
                    && \text{by \eqref{eq:C}}\\
 & = y'\circ y^*    && \text{by \eqref{eq:D}}\\
 & = y(0) y'        && \text {by Axiom 1.}
\end{align}
\end{verbatim}
A variant environment \env{alignat} allows the horizontal space between
equations to be explicitly specified. This environment takes one argument,
the number of \qq{equation columns} (the number of pairs of right-left
aligned columns; the argument is the number of pairs): count the maximum
number of \verb'&'s in any row, add~1 and divide by~2.
\begin{alignat}{2}
x& = y_1-y_2+y_3-y_5+y_8-\dots
                  &\quad& \text{by \eqref{eq:C}}\\
 & = y'\circ y^*  && \text{by \eqref{eq:D}}\\
 & = y(0) y'      && \text {by Axiom 1.}
\end{alignat}
%
\begin{verbatim}
\begin{alignat}{2}
x& = y_1-y_2+y_3-y_5+y_8-\dots
                  &\quad& \text{by \eqref{eq:C}}\\
 & = y'\circ y^*  && \text{by \eqref{eq:D}}\\
 & = y(0) y'      && \text {by Axiom 1.}
\end{alignat}
\end{verbatim}
The environment \env{flalign} (``full length alignment'') stretches the
space between the equation columns to the maximum possible width, leaving
only enough space at the margin for the equation number, if present.
\begin{flalign}
x&=y       & X&=Y\\
x'&=y'     & X'&=Y'\\
x+x'&=y+y' & X+X'&=Y+Y'
\end{flalign}
%
\begin{verbatim}
\begin{flalign}
x&=y       & X&=Y\\
x'&=y'     & X'&=Y'\\
x+x'&=y+y' & X+X'&=Y+Y'
\end{flalign}
\end{verbatim}
%
\begin{flalign*}
x&=y       & X&=Y\\
x'&=y'     & X'&=Y'\\
x+x'&=y+y' & X+X'&=Y+Y'
\end{flalign*}
%
\begin{verbatim}
\begin{flalign*}
x&=y       & X&=Y\\
x'&=y'     & X'&=Y'\\
x+x'&=y+y' & X+X'&=Y+Y'
\end{flalign*}
\end{verbatim}

%% ------------------------------------------------------------------ %%

\end{document}
